\documentclass[11pt,a4paper]{article}
\usepackage[margin=1in]{geometry}
\usepackage[T1]{fontenc}
\usepackage[utf8]{inputenc}
\usepackage{lmodern}
\usepackage{microtype}
\usepackage{enumitem}
\usepackage{xcolor}
\usepackage{hyperref}
\hypersetup{colorlinks=true, linkcolor=black, urlcolor=blue}
\setlength{\parindent}{0pt}
\setlength{\parskip}{0.55em}
\setlist[itemize]{leftmargin=1.5em,itemsep=0.2em,topsep=0.2em}
\setlist[enumerate]{leftmargin=1.7em,itemsep=0.2em,topsep=0.2em}
\pagenumbering{arabic}
\begin{document}
\begin{center}
{\LARGE\bfseries Dynamics 365 Sales CRM Migration Project Management Proposal}
\end{center}
\vspace{0.5em}

Hello,\par

My name is Fabricio Santana. I am a software engineer and technical lead with 20+ years of experience in software engineering, technical leadership, architecture, backend development, database integration, automated testing, cloud computing, and delivery of critical systems. I have also led large engineering teams and coordinated projects where business rules, data consistency, stakeholder alignment, governance, testing, deployment readiness, and operational continuity were essential.\par

Your opportunity is a strong match for my background because a CRM migration is not only a schedule-management exercise. The important part is keeping business requirements, Dynamics 365 Sales configuration, data migration, testing, go-live, hypercare, risks, issues, stakeholders, and budget under one controlled delivery model.\par

I can contribute as a hands-on Project Manager with technical depth: someone able to coordinate business stakeholders, technical teams, and external partners while also understanding the practical risks behind data migration, integrations, testing, acceptance criteria, cutover planning, and post-go-live support.\par

\section*{Working methodology}

My delivery method is designed to reduce risk before execution accelerates. For this opportunity, I recommend starting with a short Discovery Sprint / transition assessment and then confirming the right B2B engagement model based on project phase, onsite expectations, governance maturity, data complexity, and go-live pressure.\par

The first step is a one-week Discovery Sprint with a fixed price of USD 500, or an equivalent amount in EUR if preferred for local B2B contracting. In this phase, I can hold up to five remote meetings to understand the current project phase, stakeholder map, Dynamics 365 Sales scope, migration plan, data readiness, governance model, critical path, risks, open issues, budget constraints, and onsite requirements in Porto.\par

The deliverable is a practical CRM migration delivery plan, including scope assessment, stakeholder and dependency map, governance cadence, RAID structure, KPI/reporting model, testing/go-live/hypercare approach, critical risks, decision points, responsibilities, and recommended engagement model.\par

After discovery, the project can continue in one of two practical formats:\par

\begin{itemize}
\item \textbf{B2B project leadership engagement} --- part-time or full-time, depending on onsite expectations, project phase, and NextPM's delivery model. This is the best fit if the role requires continuous leadership through functional analysis, configuration, migration, testing, go-live, and hypercare.
\item \textbf{Fixed 4-week delivery-management cycle} --- a bounded engagement focused on stabilizing governance, migration planning, stakeholder alignment, reporting, RAID control, testing readiness, go-live planning, and hypercare preparation. This is useful if NextPM wants an initial controlled cycle before committing to a longer engagement.
\end{itemize}

In either model, the management approach should be hybrid and pragmatic: structured enough for PMO reporting, steering committees, budget and risk control, but flexible enough to support workshops, daily execution decisions, technical dependencies, migration changes, and fast issue resolution.\par

\section*{Opportunity analysis}

This project should be treated as a data-intensive business migration, not just a tool implementation. Dynamics 365 Sales can only deliver value if the migration preserves business continuity, data quality, sales process alignment, reporting reliability, and user adoption.\par

The visible requirement is for a Project Manager. The real delivery problem is coordinating several dependent tracks at the same time:\par

\begin{itemize}
\item business requirements and functional analysis;
\item Dynamics 365 Sales configuration and process alignment;
\item source CRM assessment, data mapping, cleansing, deduplication, migration execution, and validation;
\item integration dependencies with reporting, marketing, finance, ERP, identity, or external tools;
\item test strategy across configuration, migration data, permissions, workflows, reports, and user acceptance;
\item go-live planning, cutover responsibilities, rollback criteria, and hypercare support;
\item governance through status meetings, steering committees, workshops, KPI reporting, risk and issue control, and decision escalation.
\end{itemize}

The main risks are unclear scope boundaries, weak data ownership, delayed stakeholder decisions, underestimated migration validation, insufficient UAT, uncontrolled integration dependencies, and a go-live plan that does not define acceptance, rollback, and hypercare responsibilities clearly enough.\par

For this reason, I would start by creating a clear operating model for the migration: what decisions are needed, who owns them, what data and processes are in scope, how progress is measured, how risks are escalated, how testing is accepted, and what must be true before go-live.\par

\section*{Proposed work plan}

I would approach the engagement in structured phases. The first phase is the recommended immediate commitment; the following phases can be executed as part of the confirmed B2B engagement or fixed delivery-management cycle.\par

\textbf{1. Discovery Sprint / project transition assessment}\par

\begin{itemize}
\item Review the current project status, target go-live date, scope, stakeholders, existing plan, governance model, RAID log, budget constraints, and open decisions.
\item Clarify the Dynamics 365 Sales scope, source CRM, migration volume, data quality status, integration dependencies, and testing expectations.
\item Confirm the onsite model in Porto, remote collaboration cadence, B2B contracting expectations, allocation, and reporting responsibilities.
\item Produce a practical delivery plan with governance cadence, critical path, stakeholder map, dependency map, risk register structure, KPI/reporting model, and recommended engagement model.
\end{itemize}

\textbf{2. Governance, scope, and stakeholder alignment}\par

\begin{itemize}
\item Establish or refine project governance with status meetings, steering committees, workshops, escalation paths, decision log, and structured reporting.
\item Define scope boundaries, assumptions, dependencies, milestones, acceptance criteria, and change-control flow.
\item Align local and international stakeholders around responsibilities, delivery cadence, decision ownership, and communication channels.
\item Keep reporting focused on progress, risks, issues, dependencies, budget exposure, quality gates, and decisions required.
\end{itemize}

\textbf{3. Functional analysis and Dynamics 365 Sales delivery coordination}\par

\begin{itemize}
\item Coordinate business requirements, process mapping, functional analysis, and validation of Dynamics 365 Sales configuration decisions.
\item Ensure business requirements are translated into clear technical deliverables and acceptance criteria.
\item Track configuration dependencies, integration touchpoints, reporting expectations, permissions, and user workflows.
\item Facilitate workshops and structured follow-ups so business and technical teams remain aligned.
\end{itemize}

\textbf{4. Data migration, validation, and testing control}\par

\begin{itemize}
\item Coordinate the data migration plan, including source-data assessment, mapping, cleansing responsibilities, deduplication decisions, migration rehearsals, and validation criteria.
\item Define ownership for data quality issues, exception handling, reconciliation, and sign-off.
\item Oversee test planning across configuration, migration data, integrations, permissions, reports, UAT, and regression scenarios.
\item Track defects, blockers, retest cycles, readiness criteria, and unresolved risks before go-live.
\end{itemize}

\textbf{5. Go-live, hypercare, and handoff}\par

\begin{itemize}
\item Coordinate go-live planning, cutover checklist, readiness reviews, rollback criteria, business sign-off, support model, and communication plan.
\item Manage hypercare cadence, issue triage, escalation, ownership, prioritization, and stabilization reporting.
\item Ensure lessons learned, open items, documentation, and operational handoff are clear before closing the engagement phase.
\item Keep NextPM and the end client informed with concise, structured reporting throughout the transition.
\end{itemize}

\section*{Estimated timeline}

The immediate recommended step is a 1-week Discovery Sprint / transition assessment. After that, the engagement model should be confirmed according to the project's current phase and NextPM's onsite expectations.\par

A practical starting schedule would be:\par

\begin{itemize}
\item \textbf{Week 1:} Discovery Sprint, project transition assessment, stakeholder map, critical path, risk/dependency review, governance recommendation, and engagement proposal.
\item \textbf{Weeks 2 to 5:} first delivery-management cycle focused on governance stabilization, migration planning, data/testing readiness, reporting, stakeholder alignment, and go-live/hypercare preparation according to project phase.
\end{itemize}

If the project is already close to go-live, the first cycle should prioritize cutover readiness, data validation, UAT defects, risk escalation, and hypercare planning. If the project is earlier, the first cycle should prioritize scope control, functional analysis, configuration governance, data migration strategy, and stakeholder decision cadence.\par

\section*{Availability and communication}

I can work in Portuguese and English with business stakeholders, technical teams, and external partners. I am comfortable operating in international, cross-functional environments and keeping communication structured through workshops, status meetings, steering committees, written updates, risk/issue reports, and decision tracking.\par

One point must be confirmed early: the opportunity lists 2-3 days per week in Porto's office as a must-have requirement. The proposal should move forward only if this onsite expectation can be met or if NextPM is open to a mostly remote model with planned onsite workshops and clear travel arrangements.\par

\section*{Budget}

For the immediate step, my recommended starting budget is USD 500 for the 1-week Discovery Sprint / transition assessment, or an equivalent EUR amount if NextPM prefers local B2B invoicing. This produces a concrete CRM migration delivery plan before confirming the longer engagement model.\par

After discovery, the commercial model can be confirmed in one of these formats:\par

\begin{itemize}
\item \textbf{Fixed 4-week delivery-management cycle} from USD 4,000 when the first objective is to stabilize governance, stakeholder alignment, migration planning, reporting, testing readiness, and go-live preparation.
\item \textbf{Expanded 4-week delivery-management cycle} from USD 6,000 when the migration is complex, close to go-live, data-heavy, multi-country, integration-heavy, or requires stronger risk control and intensive stakeholder coordination.
\item \textbf{B2B day-rate or monthly engagement} if NextPM needs continuous project leadership through the full CRM migration lifecycle.
\end{itemize}

The final budget should be confirmed after understanding project duration, allocation, onsite expectations in Porto, travel arrangements, current phase, delivery pressure, and the level of responsibility expected for budget, vendor, stakeholder, and go-live control.\par

\section*{Questions before final confirmation}

Before confirming the final engagement model, I would like to clarify:\par

\begin{enumerate}
\item Is the 2-3 days per week onsite requirement in Porto mandatory every week, and are travel or local expenses covered?
\item What is the current project phase: planning, functional analysis, configuration, migration preparation, testing, go-live preparation, or recovery?
\item What is the source CRM, and what Dynamics 365 Sales scope is included in the migration?
\item What are the target go-live date and main business deadlines?
\item How many countries, business units, sales teams, or stakeholder groups are involved?
\item What data objects are in scope: accounts, contacts, leads, opportunities, activities, products, cases, custom entities, or historical records?
\item Is there already a data migration strategy, data owner, cleansing plan, and validation approach?
\item Which integrations, reports, automations, or downstream processes depend on the CRM migration?
\item Who owns functional analysis, Dynamics configuration, migration execution, testing, and final business acceptance?
\item What governance is already in place: project plan, RAID log, steering committee, KPI dashboard, decision log, and escalation path?
\item What allocation, duration, and B2B commercial model does NextPM expect for this role?
\end{enumerate}

Once these points are clear, I can confirm the recommended engagement model, milestones, reporting structure, and commercial terms.\par

\end{document}
